\scanpage{13}
\mspage{\handwritten{27}}

garantit alors que $v$ est un isomorphisme d'algébres topologiques (grâce à \deleted{\unsure{b}}\handwritten{c}) et \handwritten{d}) pour $\mathrm B$ et $\mathrm C$). Donc \deleted{les} on a 2), \inserted{donc on a 3). Or} \deleted{et} le même argument en sens inverse montre que \deleted{$u\unsure{{:}\mathrm C\to\mathrm B}$ est alors un isomorphisme d'alg top, i.e. qu'on a 4).} \inserted{3) implique 4). - Enfin 4)$\handwritten{\Rightarrow}$5) et 5)$\handwritten{\Rightarrow}$2).} On peut certainement améli\deleted{\unsure{r}}orer la rédaction de la démonstratio\deleted{\unsure{s}}n, évitant certaines redites.

Pour simplifier, nous nous restreindrons par la suite à des hypothèses noethériennes (permettant de nous déb\deleted{\unsure{aa|rr}}\unsure{ar|ra|rr|arr}asser en particulier de la désagréable hypothèse \deleted{\unsure{c}}\handwritten{d}) sur $\mathrm B$).

\uline{Définition} \deleted{\unsure{3|5}}\handwritten{4}.7. Soit $\mathrm A\handwritten{\to}\mathrm B$ un homomorphisme continu d'anneaux noethériens préadiques. On dit que $\mathrm B$ est une \uline{algèbre} (topologique) \uline{de Cohen} sur $\mathrm A$ si \inserted{\handwritten{a)}} $\mathrm B$ est séparé et complet, i.e. \inserted{\handwritten{est}} un anneau adique,\deleted{sur $\mathrm A$,} \inserted{\handwritten{b)}} s'il satisfait à la condition $\mathrm P$, et si enfin \inserted{\handwritten{c)}} $\mathrm B$ est formellement \deleted{\unsure{plat}}\inserted{simple} sur $\mathrm A$ \deleted{\unsure{On}}\inserted{\handwritten{, (ou encore, $I$ étant un \unsure{idéal} de \unsure{définition} de $A$, $B_{\unsure{0|o}}=B/IB$ est form\textsuperscript{t} simple sur $A_0=A/I$)}}

On notera qu'en vertu de 3.1. , si $\mathrm I$ est un idéal de définition de $\mathrm A$, la dernière condition équivaut à la condition que $\mathrm B/\mathrm{IB}$ est formellement simple sur $\mathrm A/\mathrm{IA}$. D'ailleurs, on peut interpréter \deleted{3}\handwritten{4}.6. de la façon suivante (et il est \deleted{peutêtre} plus raisonnable de réléguer \deleted{3}\handwritten{4}.6. en lemme, et de réserver le corollaire en théorème):

\uline{Coroll\deleted{\unsure{e}}aire} \deleted{3}\handwritten{4}.8. Soient $\mathrm A$ un anneau \deleted{préadique} noethérien, $\mathrm I$ un idéal de $\mathrm A$, munissons $\mathrm A$ de la topologie $\mathrm I$-préadique, posons $\mathrm A_0=\mathrm A/\mathrm{IA}$, et soit $\mathrm B_0$ une $\mathrm A_0$-algèbre topologique. \deleted{S'il}\inserted{Supposons qu'\handwritten{il}} existe une algèbre de Cohen $\mathrm B$ sur $\mathrm A$ telle que $\mathrm B/\mathrm{IB}$ soit isomorphe à $\mathrm B_0$ comme $\mathrm A_0$-algèbre topologique,\inserted{(ce qui implique $\mathrm B_0$ \deleted{\unsure{formt simple}}\handwritten{anneau de Cohen} sur $\mathrm A_0$).} alors $\mathrm B$ est unique à \inserted{un} isomorphisme près \deleted{(induisant l'isomorphisme donné sur} (compatible avec les isomorphismes donnés des algèbres \deleted{\unsure{réduite|réduites}} réduites mod $\mathrm I$ avec $\mathrm B_0$). De plus, si $\mathrm B$ est une $\mathrm A$-algèbre topologique adique telle que $\mathrm B/\mathrm{IB}$ soit \deleted{\unsure{A}}isomorphisme à $\mathrm B_0$ comme $\mathrm A_0$-algèbre topologique, alors les conditions suivantes sur $\mathrm B$ sont équivalentes: \handwritten{\unsure{0|a})} $\mathrm B$ est une $\mathrm A$-algèbre de Cohen
% archive sheet number: 12
